Sei
$\Sigma = \{\texttt{a},\texttt{b}\}$
und
$\Sigma_1 = \{\texttt{a},\texttt{b},\texttt{\#}\}= \Sigma\cup\{\texttt{\#}\}$.
Ist die Sprache
\[
L
=
\{
v\texttt{\#}t
\mid
\text{$t$ ist Teilzeichenkette von $v$}
\}
\]
kontextfrei?

\begin{hinweis}
Betrachten Sie Wörter der Form
$\texttt{a}\mathstrut^n\texttt{b}\mathstrut^n\texttt{\#}\texttt{a}\mathstrut^n\texttt{b}\mathstrut^n$.
\end{hinweis}

\begin{loesung}
\definecolor{darkred}{rgb}{0.8,0,0}
\definecolor{darkgreen}{rgb}{0,0.6,0}
Die Sprache ist nicht kontextfrei, wie man mit dem Pumping Lemma für
kontextfreie Sprachen beweisen kann.
\def\h{0.6}
\def\unterteilung#1{
	\fill[color=blue!20,opacity=0.7]
		({0*\h+0.05},0.05) rectangle ({(#1)*\h-0.025},{\h-0.05});
	\draw[color=blue]
		({0*\h+0.05},0.05) rectangle ({(#1)*\h-0.025},{\h-0.05});
	\node[color=blue] at ({0.5*(#1)*\h},{0.5*\h}) {$u\mathstrut$};

	\fill[color=darkred!20,opacity=0.7]
		({(#1)*\h+0.025},0.05) rectangle ({(#1+1)*\h-0.025},{\h-0.05});
	\draw[color=darkred]
		({(#1)*\h+0.025},0.05) rectangle ({(#1+1)*\h-0.025},{\h-0.05});
	\node[color=darkred] at ({(#1+0.5)*\h},{0.5*\h}) {$v\mathstrut$};

	\fill[color=darkgreen!20,opacity=0.7]
		({(#1+1)*\h+0.025},0.05)
			rectangle ({(#1+2.5)*\h-0.025},{\h-0.05});
	\draw[color=darkgreen]
		({(#1+1)*\h+0.025},0.05)
			rectangle ({(#1+2.5)*\h-0.025},{\h-0.05});
	\node[color=darkgreen] at ({(#1+1.75)*\h},{0.5*\h}) {$x\mathstrut$};

	\fill[color=darkred!20,opacity=0.7]
		({(#1+2.5)*\h+0.025},0.05)
			rectangle ({(#1+3.5)*\h-0.025},{\h-0.05});
	\draw[color=darkred]
		({(#1+2.5)*\h+0.025},0.05)
			rectangle ({(#1+3.5)*\h-0.025},{\h-0.05});
	\node[color=darkred] at ({(#1+3)*\h},{0.5*\h}) {$y\mathstrut$};

	\fill[color=blue!20,opacity=0.7]
		({(#1+3.5)*\h+0.025},0.05) rectangle ({21*\h-0.05},{\h-0.05});
	\draw[color=blue]
		({(#1+3.5)*\h+0.025},0.05) rectangle ({21*\h-0.05},{\h-0.05});
	\node[color=blue] at ({0.5*(#1+24.5)*\h},{0.5*\h}) {$u\mathstrut$};

	\draw[line width=0.2pt] ({(#1)*\h},0) -- ++(0,-0.3);
	\draw[line width=0.2pt] ({(#1+3)*\h},0) -- ++(0,-0.3);
	\draw[<->] ({(#1)*\h},-0.2) -- ({(#1+3)*\h},-0.2);
	\node at ({(#1+1.5)*\h},-0.2) [below] {$\le N\mathstrut$};
}
\def\hintergrund{
	\draw (0,0) rectangle ({21*\h},\h);
	\draw ({5*\h},0) -- ++(0,\h);
	\draw ({10*\h},0) -- ++(0,\h);
	\draw ({11*\h},0) -- ++(0,\h);
	\draw ({16*\h},0) -- ++(0,\h);
	\node at ({2.5*\h},{0.5*\h}) {\texttt{a}\strut};
	\node at ({7.5*\h},{0.5*\h}) {\texttt{b}\strut};
	\node at ({10.5*\h},{0.5*\h}) {\texttt{\#}\strut};
	\node at ({13.5*\h},{0.5*\h}) {\texttt{a}\strut};
	\node at ({18.5*\h},{0.5*\h}) {\texttt{b}\strut};
	\node at ({5*\h},{\h-0.1}) [above] {$N\mathstrut$};
	\node at ({10.5*\h},{\h-0.1}) [above left] {$2N\mathstrut$};
	\node at ({10.5*\h},{\h-0.1}) [above right] {$2N+1\mathstrut$};
	\node at ({16*\h},{\h-0.1}) [above] {$3N+1\mathstrut$};
	\node at ({21*\h},{\h-0.1}) [above] {$4N+1\mathstrut$};
}
\begin{enumerate}
\item
Annahme: $L$ ist kontextfrei
\item
Nach dem Pumping Lemma gibt es die Pumping Length $N$
\item
Wir wählen das Wort
$\texttt{a}\mathstrut^N\texttt{b}\mathstrut^N\texttt{\#}\texttt{a}\mathstrut^N\texttt{b}\mathstrut^N$.
\begin{center}
\begin{tikzpicture}[>=latex,thick]
\hintergrund
\end{tikzpicture}
\end{center}
\item
Nach dem Pumping Lemma gibt es eine Aufteilung
\(
w=
{\color{blue}u}
{\color{darkred}v}
{\color{darkgreen}x}
{\color{darkred}y}
{\color{blue}z}
\)
mit
$| {\color{darkred}v} {\color{darkgreen}x} {\color{darkred}y} |\le N$
und
$| {\color{darkred}v} {\color{darkred}y} |>0$
\item
Es gibt drei Fälle:
\begin{enumerate}[a)]
\item
Die roten Teile liegen beide vor dem \texttt{\#}-Zeichen.
\begin{center}
\begin{tikzpicture}[>=latex,thick]
\hintergrund
\unterteilung{1.75}
\end{tikzpicture}
\end{center}
Abpumpen führt dazu, dass der Teil links vom \texttt{\#} jetzt kürzer ist
und damit den Teil rechts davon nicht mehr als Teilstring enthalten kann.
\item
Die roten Teile liegen beide nach dem \texttt{\#}-Zeichen.
\begin{center}
\begin{tikzpicture}[>=latex,thick]
\hintergrund
\unterteilung{14.75}
\end{tikzpicture}
\end{center}
Aufpumpten führt dazu, dass der Teil rechts vom \texttt{\#} nun grösser
ist, also nicht mehr Teilstring sein kann.
\item
Die beiden Teile roten Teile liegen im Block von \texttt{b} vor dem
\texttt{\#} und/oder im Block von \texttt{a} nach dem \texttt{\#}.
\begin{center}
\begin{tikzpicture}[>=latex,thick]
\hintergrund
\unterteilung{8.75}
\end{tikzpicture}
\end{center}
Durch Abpumpen wird die Anzahl der \texttt{b} links vom \texttt{\#}
kleiner und kann daher das Wort rechts von \texttt{\#} nicht mehr erhalten.
Durch Aufpumpen wird die Anzahl der \texttt{a} rechts vom \texttt{\#}
grösser und kann daher nicht mehr im Wort links von \texttt{\#} enthalten
sein.
\end{enumerate}
Das gepumpte Wort ist also nicht mehr in $L$ im Widerspruch zur
Aussage des Pumping-Lemma.
\item
Dieser Widerspruch zeigt, dass $L$ nicht kontextfrei sein kann.
\qedhere
\end{enumerate}
\end{loesung}

\begin{bewertung}
Annahme $L$ kontextfrei ({\bf PL}) 1 Punkt,
Pumping Length ({\bf N}) 1 Punkt,
geeignetes Wort wählen ({\bf W}) 1 Punkt,
Unterteilung des Wortes ({\bf U}) 1 Punkt,
Widerspruch beim Pumpen ({\bf P}) 1 Punkt, dieser Punkte wurde nicht
gegeben, wenn die Teilzeichenketteneigenschaft nicht untersucht wurde,
Schlussfolgerung ({\bf S}) 1 Punkt.
\end{bewertung}

